\section{Related Work}

Our work builds on three research threads: data attribution methods, attribution benchmarking, and the analysis of learning dynamics.

\subsection{Data Attribution Methods}

Training data attribution has evolved from influence functions \cite{koh2017understanding} through increasingly scalable approximations. TracIn \cite{pruthi2020estimating} introduced gradient tracing across checkpoints, trading theoretical guarantees for computational tractability. TRAK \cite{park2023trak} achieved further speedup through random projection of gradients, using the insight that influence rankings may be preserved even when absolute scores are approximate. Kronfluence \cite{grosse2023studying} applied Kronecker-factored curvature approximation (K-FAC) to enable influence computation at LLM scale.

These methods represent different computational strategies, but prior work has not systematically characterized \emph{what} each strategy measures. Our work addresses this gap by showing that computational differences manifest as systematic differences in sensitivity to influence modes.

\subsection{Attribution Method Fragility and Benchmarking}

Basu et al.\ \cite{basu2020influence} demonstrated that influence functions are fragile in deep networks, with approximation errors increasing with network depth. This finding raised concerns about the reliability of gradient-based attribution. However, we show that apparent fragility may be reframed: rather than random errors, different methods exhibit \emph{systematic} sensitivity patterns that become predictable fingerprints.

The DATE-LM benchmark \cite{jiao2025datelm} provided the first unified evaluation of LLM attribution methods, revealing that no single method dominates across tasks. Our work complements DATE-LM by explaining \emph{why} methods differ: the inductive biases embedded in each algorithm's mathematics create different sensitivities to memorization, feature transfer, and spurious association.

Several efficiency-focused works have pushed attribution to larger scales. LoRIF \cite{li2026lorif} achieved 20x storage reduction on 70B models, and GraSS \cite{hu2025grass} demonstrated 165\% throughput improvement through gradient sparsification. GGDA \cite{ley2024ggda} introduced group-level attribution with 10--50x speedup. These works focus on computational efficiency; our work focuses on characterizing what efficient methods actually measure.

\subsection{Learning Dynamics and Influence Modes}

Research on learning dynamics has identified distinct phases and patterns in neural network training. Studies of memorization \cite{feldman2020does} show that models memorize atypical examples to achieve low training loss. Feature transfer research demonstrates how learned representations enable generalization. Work on spurious correlations \cite{sagawa2020distributionally} reveals that models exploit dataset shortcuts.

These distinct influence \emph{modes}---memorization, feature transfer, spurious association---have been studied in isolation. Our contrastive mode probing framework provides a unified lens for measuring how attribution methods respond to each mode, enabling systematic comparison.

\subsection{Positioning of Our Work}

Prior work either develops individual attribution methods (TRAK, TracIn, Kronfluence), benchmarks their accuracy (DATE-LM), or addresses their efficiency (LoRIF, GraSS). We provide the missing piece: a framework for characterizing methods by their sensitivity profiles, explaining why benchmarks show method-dependent performance and enabling informed method selection based on application requirements.
